\documentclass[12pt]{article}

\usepackage{latexsym}
\usepackage[T1]{fontenc} 
\usepackage[utf8]{inputenc} 
\usepackage[spanish]{babel}
\usepackage{amsthm, amssymb, amsmath, amsfonts, bbm, mathtools} 
\usepackage{multirow, array}
\usepackage{stmaryrd}
\usepackage{xcolor}
\decimalpoint

\usepackage{wrapfig}
\usepackage{subcaption}

% no tocar nada de aqui
\providecommand{\abs}[1]{\lvert#1\rvert}
\usepackage[hidelinks]{hyperref}
\usepackage{enumerate} 
\usepackage{graphicx}

% entornos personalizados
\newtheorem{example}{Ejemplo}
\newtheorem{notacion}{Notación}
\newtheorem{defi}{Definición}
\newtheorem{obse}{Observación}
\newtheorem{lem}{Lema}

\newcommand{\sol}{\textit{Solución. }}
\renewcommand{\qed}{\hfill $\blacksquare$}
\newcommand{\eej}{\hfill $\square$}

%--------------------- Comandos personalizados ---------------------

% Simbolos de los Numeros
\newcommand{\D}{\mathfrak{D}}
\newcommand{\A}{\mathcal{A}}
\newcommand{\C}{\mathbb{C}}
\newcommand{\F}{\mathbb{F}}
\newcommand{\I}{\mathbb{I}}
\newcommand{\N}{\mathbb{N}}
\newcommand{\Q}{\mathbb{Q}}
\newcommand{\R}{\mathbb{R}}
\newcommand{\Z}{\mathbb{Z}}

% Signos de agrupacion
\newcommand{\pp}[1]{\left( #1 \right)}
\newcommand{\cc}[1]{\left[ #1 \right]}
\newcommand{\set}[1]{\left\{#1\right\}}

\newcommand{\ol}[1]{\overline{#1}}
\newcommand{\ceil}[1]{\left\lceil#1\right\rceil}
\newcommand{\floor}[1]{\left\lfloor#1\right\rfloor}
\newcommand{\ind}[1]{\mathbbm{1}_{\{#1\}}}
\newcommand{\norm}[1]{\left\Vert #1 \right\Vert}

\newcommand{\Text}[1]{\quad\text{#1}\quad}

\newcommand{\bs}{\backslash}
\renewcommand{\c}{\complement}
\newcommand{\beq}{bilipschitz }

% Flechas
\newcommand{\imp}{\Longrightarrow}
\newcommand{\Imp}{\quad\Longrightarrow\quad}
\newcommand{\ssi}{\Longleftrightarrow}
\newcommand{\Ssi}{\quad\Longleftrightarrow\quad}


\newcommand{\e}{\varepsilon}
\renewcommand{\d}{\delta}
\renewcommand{\O}{\Omega}
\renewcommand{\l}{\lambda}
\renewcommand{\a}{\alpha}
\renewcommand{\b}{\beta}
\renewcommand{\abs}[1]{\left\lvert#1\right\rvert}
%\renewcommand{\qedsymbol}{$\blacksquare$}


\newcommand{\Co}{\mathfrak C}
\newcommand{\T}{\mathbb T}
%\newcommand{\uno}{\mathbb 1}
%\newcommand{\ind}[1]{\mathbbm{1}_{\{#1\}}}
\newcommand{\uno}{\mathbbm{1}}


\renewcommand{\o}{\omega}

\newcommand{\x}{\mathrm x}
\newcommand{\y}{\mathrm y}
\newcommand{\Int}{\mathrm{int}}
\newcommand{\Ext}{\mathrm{ext}}
\newcommand{\diam}{\mathrm{diam}}
\newcommand{\cl}{\mathrm{cl}}
\newcommand{\angles}[2]{\left\langle #1, #2\right\rangle}

\textheight=25cm \textwidth=18cm \topmargin=-2cm \oddsidemargin=-1cm

\setlength{\parskip}{0.8em}  % Espaciado párrafos
\setlength{\parindent}{0pt}  % Sangría

\newcommand{\To}{\to}



\begin{document}
    \pagestyle{empty}
    \begin{center}
        {\textsc{Ecuaciones Diferenciales Parciales II} \\
        \medskip 
        \large{Tarea 3}}
    \end{center}
    
    \bigskip
    
    \begin{enumerate}
        \item Resolver en Python por el método de Diferencias finitas los
        siguientes ejercicios, de preferencia métodos implíctos:
        \begin{itemize}
            \item Ecuación con condiciones de Dirichlet,
            \begin{align*}
                \frac{\partial u}{\partial t} &= 1
                *\frac{\partial^2u}{\partial x^2},\\
                x\in (0,1),&\quad 0,2\geq t > 0,\\
                u(x,0) &= 4x+4x^2,\quad x\in[0,1],\\
                U(0, t) &= 0,\quad 0,2 \geq t \geq 0,\\
                U(1, t) &= 0,\quad 0,2 \geq t \geq 0.
            \end{align*}

            \item Ecuación con condiciones de Neumann,
            \begin{align*}
                \frac{\partial u}{\partial t} &= 1
                *\frac{\partial^2u}{\partial x^2},\\
                x\in (0,1),&\quad 0,2\geq t > 0,\\
                u(x,0) &= 4x+4x^2,\quad x\in[0,1],\\
                U_x(0, t) &= 0,\quad 0,2 \geq t \geq 0,\\
                U_x(1, t) &= 0,\quad 0,2 \geq t \geq 0.
            \end{align*}

            \item Ecuación con condiciones de frontera periódicas,
            \begin{align*}
                \frac{\partial u}{\partial t} &= 1
                *\frac{\partial^2u}{\partial x^2},\\
                x\in (0,1),&\quad 0,15\geq t > 0,\\
                u(x,0) &= 4x+4x^2,\quad x\in[0,1],\\
                U(0, t) &= u(1,t),\quad 0,15 \geq t \geq 0,\\
                U_x(0, t) &= u_x(1,t),\quad 0,15 \geq t \geq 0.
            \end{align*}
        \end{itemize}
        
    \item Del libro Solving PDEs in Python: The Fenics Tutorial 1, elije dos
    ejemplos y reprodúcelos, explica el procedimiento, puedes modificar los
    ejemplos si consideras adecuado, como la malla, constantes, etc.
    \end{enumerate}

    \sol Solicones implementadas
    \href{https://github.com/HazelS002/EDP-II/tree/main}{Aquí}\eej
\end{document}
